% Generated draft: requires mathematical and source review.
\documentclass[11pt]{article}
\usepackage[margin=25mm]{geometry}
\usepackage{fontspec}
\setmainfont{Latin Modern Roman}
\usepackage{amsmath,amssymb,mathtools,unicode-math}
\setmathfont{Latin Modern Math}
\usepackage{enumitem,needspace,xurl,hyperref}
\setlength{\parindent}{0pt}
\setlength{\parskip}{5pt}
\setlength{\emergencystretch}{3em}
\newcommand{\N}{\mathbb{N}}
\newcommand{\Z}{\mathbb{Z}}
\newcommand{\Q}{\mathbb{Q}}
\newcommand{\R}{\mathbb{R}}
\newcommand{\C}{\mathbb{C}}
\newcommand{\F}{\mathbb{F}}
\newcommand{\A}{\mathbb{A}}
\newcommand{\PP}{\mathbb{P}}
\newcommand{\E}{\mathbb{E}}
\newcommand{\eps}{\varepsilon}
\newcommand{\dd}{\,\mathrm{d}}
\begin{document}
\section*{UnsolvedMath ID: 30002400}
\textbf{Source problem code:} catalogue entry
\subsection{Definitions and mathematical statement}
A graded matrix factorization of a homogeneous potential \(w\) over a graded ring \(R\) consists of graded modules and maps whose two compositions are multiplication by \(w\), with grading shifts chosen to make these identities homogeneous. The supplied record does not specify \(R\), the base field, the potential, the grading convention, or the construction of the category of factorizations.

For clarity, the following is a conditional framework for interpreting the entropy notation, rather than additional hypotheses attributed to the source. Let \(\mathcal C\) be a triangulated category with shift \([1]\) and split generator \(G\), meaning that the smallest full triangulated subcategory containing \(G\) and closed under direct summands is \(\mathcal C\). For \(t\in\R\), define
\[
\delta_t(G,E)=\inf\left\{\sum_{j=1}^{m}e^{t n_j}:\begin{array}{l}0=E_0,E_1,\ldots,E_m,\quad E_m\cong E\oplus E',\\ E_{j-1}\longrightarrow E_j\longrightarrow G[n_j]\longrightarrow E_{j-1}[1]\\ \text{are distinguished triangles},\quad n_j\in\Z\end{array}\right\}.
\]
Here the infimum runs over finite towers and auxiliary objects \(E'\); an empty admissible set has infimum \(+\infty\). For an exact endofunctor \(F\colon\mathcal C\to\mathcal C\), a usual categorical entropy convention is
\[
h_t(F)=\lim_{n\to\infty}\frac{1}{n}\log\delta_t(G,F^nG),
\]
provided the limit and the required independence of \(G\) hold in the chosen framework. No existence, finiteness, or independence assertion is supplied here. If \(\mathcal C\) is linear over a field \(k\) and has finite-dimensional morphism spaces, a Serre functor is an autoequivalence \(S\) equipped with bifunctorial isomorphisms
\[
\operatorname{Hom}_{\mathcal C}(X,Y)^*\cong\operatorname{Hom}_{\mathcal C}(Y,SX),\qquad (-)^*=\operatorname{Hom}_k(-,k).
\]
The source does not state hypotheses ensuring the existence of such a functor or the applicability of this entropy convention.

\textbf{Source question, with unresolved scope retained.} For general categories \(\mathcal C\) of graded matrix factorizations for which the relevant Serre functor and categorical entropy are defined, determine a geometric interpretation of
\[
t\longmapsto h_t(S).
\]
More generally, determine whether \(\exp(h_0(F))\) is always algebraic and identify the conditions under which algebraicity holds. If \(h_0(F)\) is finite and real, the conventional interpretation of algebraicity is
\[
\exists\,P\in\Q[x]\setminus\{0\}\quad P\bigl(\exp(h_0(F))\bigr)=0.
\]
The source does not specify the admissible categories or functors \(F\) for this broader question, the intended meaning of a geometric interpretation, or conventions for infinite entropy. Thus this is an explicitly qualified research question, not a fully specified universal conjecture.
\subsection{Short English statement}
The problem asks for a geometric meaning of the categorical entropy of Serre functors on categories of graded matrix factorizations, and asks when the exponential of entropy at \(t=0\) is algebraic. The precise categories, functors, entropy hypotheses, and intended geometric interpretation remain unspecified. The supplied record labels the problem open; this draft does not independently verify that status.
\subsection{Sources}
\begin{itemize}
\item Supplied local UnsolvedMath registry record: permanent ID \(30002400\), title \emph{Entropy of Serre Functors and Algebraicity}, source problem code \texttt{catalogue entry}, in the \emph{Oberwolfach Reports Open Problems} collection. The record gives proposed year 2013 and names no individual proposer.
\item Supplied report citation: \emph{Matrix Factorizations in Algebra, Geometry, and Physics} (2014), \emph{Oberwolfach Reports}, pp.\ 2501--2552, DOI \texttt{10.4171/owr/2013/44}. \url{https://doi.org/10.4171/owr/2013/44}. The supplied background describes these reports as documenting mathematical workshops and their open-problem discussions. No location within the report is supplied for this particular question.
\item Supplied literature source: \emph{Oberwolfach Report}. \url{https://ems.press/content/serial-article-files/46474}.
\item Supplied literature source: \emph{Higher hereditary algebras}. \url{https://arxiv.org/abs/2201.00302}. The record supplies no specific theorem or passage connecting this reference to the question.
\item Literature assessment dated exactly \texttt{2026-08-21}: the supplied curation labels the problem open and states that special computations relate entropy of Serre functors to categorical growth, while no general geometric interpretation for the broader class of categories in the extracted question is established. This is supplied dated curation derived from the OWR material, not an independently verified assessment; its time-sensitive status and cited claims require independent verification.
\item The registry assigns statement status \texttt{corrected\_verified} and reports that the mathematical content was checked against the official report, with notation, context, and question form normalized without changing the source-level claim. This is a supplied verification assertion, not verification performed for this draft.
\end{itemize}
\end{document}
